% 09b_wayland
\section{图形合成器栈：Weston 于 StarryOS}

Wayland 是取代 X 窗口系统的显示服务器协议，其中的合成器是介于内核与各图形应用之间、独占显示硬件的进程：它从各客户端收集绘制好的窗口缓冲，叠合为一帧整屏画面交给显示设备，同时把键盘、鼠标与触摸输入分发回对应窗口。让一个完整的合成器在一套新内核上启动并接入真实客户端，等于一次性把 Linux 图形与进程间通信栈的一大部分纳入内核的支撑范围：它对内核的要求横跨显示输出、输入设备发现、套接字传递文件描述符与事件循环四个彼此独立的子系统，任一处不合规整套界面即无法正常运行。StarryOS 此前不具备承载这类合成器所需的内核接口，图形应用因而无从运行。现已\ours{在 StarryOS 上完整启用 Weston 14 图形合成器}，Weston 是 Wayland 的参考实现，它以 DRM 加 pixman 软件后端启动、接受客户端连接、经 GTK4 与 Xwayland（xeyes）客户端呈现界面并支持 VNC 远程查看，并在四种体系结构上通过自动化判定。本节先结合图~\ref{fig:wayland} 说明支撑合成器的四个内核子系统及其机制，再说明逐层的内核启用与修复，最后给出验证。

\figphl{../figures/out/Fwayland.png}{Weston 于 StarryOS：客户端经 \texttt{AF\_UNIX} 套接字连入并提交窗口缓冲，输入沿 libinput 与 libudev 自 sysfs 与 evdev 发现，DRM/KMS 承载扫描输出，整套事件经 \texttt{epoll} 与 \texttt{timerfd} 汇聚轮转}{fig:wayland}

\subsection{支撑合成器的四个内核子系统}

这四个子系统各自对应一组内核接口，以下逐一说明其机制。

\paragraph{显示输出 DRM/KMS}
显示后端 \texttt{drm-backend} 经 libdrm 操作 \texttt{/dev/dri/card0}。它要协商设备能力、枚举 KMS（内核态显示模式设置）资源、分配显示缓冲，并以原子方式提交一次模式设置与翻页。模式设置须兼容传统的 \texttt{SETCRTC} 路径与原子化的 KMS 两种语义，承载像素的显示缓冲则由 GEM dumb 缓冲提供，合成器把叠合好的整屏画面写入这块缓冲后交由显示控制器扫出。

\paragraph{输入发现 evdev 与 libinput}
输入由 libinput 经 libudev 沿 sysfs 遍历发现 \texttt{/dev/input/eventN} 节点，再以一组 evdev ioctl 查询各设备能力并读取事件。其中 \texttt{EVIOCGPROP} 用于取回设备属性、\texttt{EVIOCGABS} 用于取回绝对轴信息，二者是 libinput 判定一个节点是否为可用输入设备、以及该设备属键盘、指点还是触摸的依据；缺此则该节点被静默跳过。

\paragraph{协议传输 AF\_UNIX 与文件描述符传递}
协议服务由 wayland-server 在一个 \texttt{AF\_UNIX} 流式套接字上监听，接受客户端连接并接收其提交的窗口缓冲。其中最关键的是随消息原子地传递文件描述符：客户端把承载像素的共享内存以文件描述符的形式经辅助控制消息交给合成器，套接字须在字节流中以正确的标记原子地携带这一控制消息，缺此则合成器无从取得客户端的窗口缓冲。共享内存一侧还依赖 memfd 的密封语义，\texttt{F\_SEAL\_SHRINK}、\texttt{F\_SEAL\_GROW} 与 \texttt{F\_SEAL\_WRITE} 分别封住缩小、增长与写入，令合成器得以安全映射客户端提交的缓冲而不虞其被抽换。

\paragraph{事件循环 epoll 与 timerfd}
合成器靠 \texttt{epoll} 与 \texttt{timerfd} 把上述三路事件汇入一个统一的事件循环轮转。\texttt{epoll} 多路复用显示、输入与套接字三类文件描述符的就绪状态，\texttt{timerfd} 则把定时器表达为一个可被 \texttt{epoll} 等待的文件描述符，用于驱动重绘节拍与各类超时。这一层要求就绪通知与实际就绪的数据严格一致，否则事件循环会在虚假就绪上空转。

\subsection{逐层的内核启用与修复}

此次适配遇到的问题落在三个深浅不同的层次，内核改动据此逐层推进。最浅的一层是接口不存在，打开即以明确错误码回绝。系统内核侧补上了 \texttt{/dev/dri/card0} 节点，兼容 \texttt{SETCRTC} 与原子 KMS 两条模式设置路径（\href{https://github.com/rcore-os/tgoskits/pull/506}{\#506}），并以按缓冲粒度分配的 GEM dumb 缓冲与带引用计数的 mmap 承载像素（\href{https://github.com/rcore-os/tgoskits/pull/514}{\#514}）；输入侧如实暴露多个 \texttt{eventN} 设备并实现 \texttt{EVIOCGPROP} 与 \texttt{EVIOCGABS}（\href{https://github.com/rcore-os/tgoskits/pull/513}{\#513}），设备发现所依赖的 \texttt{AF\_NETLINK} 上 rtnetlink、uevent 与 genl 几路也一并补全（\href{https://github.com/rcore-os/tgoskits/pull/512}{\#512}），以上均已合并。

第二层是接口看似齐备、编译与打开均成功而行为错误且不报任何错。sysfs 的 \texttt{/sys/\allowbreak class/\allowbreak input/} 须是指向真实设备路径的符号链接，libudev 靠 \texttt{realpath()} 与 \texttt{dirname()} 沿链接回溯，若做成普通目录回溯即中途断掉；\texttt{/dev/\allowbreak input/\allowbreak eventN} 的次设备号须从 \texttt{EVDEV\_MINOR\_BASE}（64）起编，若从 1 起编，libinput 以 \texttt{fstat()} 的 \texttt{st\_rdev} 反查 \texttt{/sys/\allowbreak dev/\allowbreak char/} 便对不上号。\ours{这些隐式约定集中在一处修订内补齐}：sysfs 类目录改为符号链接、evdev 次设备号改从 64 起编、并预先在 \texttt{/run/udev} 下写入 udev 的初始化种子文件，libudev 方肯认定设备已配置（\href{https://github.com/rcore-os/tgoskits/pull/508}{\#508}，已合并）。memfd 的密封检查被分布到截断、映射与写入各条路径，缺一处密封即形同虚设（\href{https://github.com/rcore-os/tgoskits/pull/507}{\#507} 与 \href{https://github.com/rcore-os/tgoskits/pull/515}{\#515}，已合并）。Weston 的整体启动收尾、把输入接到中断唤醒以降低延迟、以及让 \texttt{AF\_UNIX} 流式套接字按字节标记原子地携带控制消息送达文件描述符，同在一处合并完成（\href{https://github.com/rcore-os/tgoskits/pull/509}{\#509}，已合并）。这一层的失败是沉默的，libinput 只报一句跳过未配置的输入设备，要查出真因须转而阅读用户态库的源码。

最深的一层是并发与性能，它不挡启动却使界面难以使用。\texttt{epoll} 的水平触发注册无条件唤醒等待方，产生每秒约 500 次的空转忙循环；输入若退回内核节拍驱动会带来约 16 毫秒的输入延迟，在单核上还叠加中断上下文中的锁争用。系统内核侧把 \texttt{epoll} 的水平触发改为如实排干就绪事件，收敛了那条 500 赫兹的空转（\href{https://github.com/rcore-os/tgoskits/pull/504}{\#504}），并补上 \texttt{timerfd} 这一计时事件源（\href{https://github.com/rcore-os/tgoskits/pull/503}{\#503}），二者均已合并。输入延迟已由中断唤醒降低（\href{https://github.com/rcore-os/tgoskits/pull/509}{\#509}），单核上中断上下文中的锁争用仍是有待优化的残留限制。

\subsection{验证}

这套支撑经端到端验证为可用。Weston 以 DRM 加 pixman 后端启动后，一个客户端可成功连接，一个 GTK4 程序能把界面绘制出来并经 VNC 远程查看，一套自动化用例产出 \texttt{WAYLAND\_TEST\_PASSED} 判据，\keyres{在四种体系结构上一致通过}（\href{https://github.com/rcore-os/tgoskits/pull/1160}{\#1160}，已合并）。为防回归，另引入以黄金帧作视觉比对的持续集成哨兵守护该能力不回归。Xwayland 亦成功运行 xeyes 这个经典 X 程序，使既有的 X 应用得以在 Wayland 之上运行（\href{https://github.com/rcore-os/tgoskits/pull/516}{\#516}，已合并）。该图形栈的启用与视觉感知流水线相互独立，为设备侧的实时监控与可视化界面提供了用户侧通路。
